\section{Route evaluation and next theorem}

Theorem \ref{thm:germ} is a genuine all-orbit analytic advance.  It solves
the varying-field topology problem left open by the finite tagged ledger and
makes rational norm pushforward continuous.  Theorem \ref{thm:boundary}
simultaneously explains why the resulting object is a two-variable germ
rather than an ungraded zeta function.

The strict Route-A tuple is
\[
(\mathrm{A1\_WEAK},\mathrm{A2\_FAIL},
\mathrm{A3\_PARTIAL\_ANALYTIC\_STRUCTURE},
\mathrm{A4\_FORMAL\_HINT}),
\]
with overall verdict \textbf{ROUTE\_A\_EXPLORATORY}.  The entries have the
following meanings.
\begin{itemize}
\item A1 is weak: the packet clock is intrinsic, but there is no
  all-rational-prime orbit correspondence.
\item A2 fails: a Banach-valued series is not a transfer operator or
  Fredholm determinant.
\item A3 is partial: normal convergence and joint holomorphy are proved in
  a right-half-plane times the unit disk, but no continuation, functional
  equation, zero law, or von-Mangoldt trace follows.
\item A4 remains a formal hint: exact $SL_2$ monodromy and residue orders do
  not define a Hilbert--P\'olya operator.
\end{itemize}

The next theorem should address the Abel boundary rather than generate more
finite packet rows.  Natural possibilities include a source-native
subtraction, a logarithmic derivative, a distributional boundary value, or
a Tauberian normalization as $u\uparrow1$.  Any acceptable construction
must retain the orbit, index, and prime-ideal tags until after the boundary
operation and must connect to the pressure/von-Mangoldt mass scale without
fitting target prime data.

\paragraph{Conclusion.}
The finite packet ledger has become an all-orbit analytic germ.  The first
remaining wall is now equally precise: the raw Abel boundary diverges on one
certified orbit.  This positive theorem plus scoped obstruction identifies a
single, testable next bridge.
